\documentclass[11pt,a4paper,UTF8]{ctexart}
\usepackage{amsmath,amssymb,mathtools,geometry,hyperref,enumitem,xcolor}
\geometry{margin=2cm}
\hypersetup{hidelinks}
\setlength{\parindent}{0pt}
\setlength{\parskip}{6pt}
\linespread{1.25}
\newcommand{\R}{\mathbb R}
\newcommand{\eps}{\varepsilon}
\newcommand{\idea}[1]{\par\smallskip\noindent\colorbox{yellow!22}{\textbf{一句话记住：}#1}\par\smallskip}
\newcommand{\warn}[1]{\par\smallskip\noindent\textcolor{red!70!black}{\textbf{坑：}#1}\par\smallskip}
\title{数学分析（上册）四章\quad 零基础也能懂\\[4pt]\large 从“数”讲到“导数”}
\author{给完全没学过的人看的}
\date{}
\begin{document}
\maketitle
\tableofcontents\newpage

\section{开场：这本书在讲什么}

数学分析，说白了就是研究两件事：\textbf{“东西跑到哪去”}（极限）和\textbf{“变化有多快”}（导数）。

你可以把它想成学开车：先认识油门、刹车（数的基本性质），再学会看仪表（极限），最后学会踩油门加速的感觉（导数）。

前四章的顺序其实就是一条线：
\begin{enumerate}[leftmargin=1.4em]
\item 第 1 章 \textbf{集合与映射}——先把“东西”和“东西之间的对应关系”说清楚。这是盖房子的砖。
\item 第 2 章 \textbf{数列极限}——一堆数排成队，它们最后趋向哪里？（数的“跑步终点”）
\item 第 3 章 \textbf{函数极限与连续}——曲线走啊走，最后停在哪、会不会断？（函数的“跑步终点”）
\item 第 4 章 \textbf{微分}——曲线在每一点有多陡？油门踩多狠？（斜率与变化率）
\end{enumerate}

\idea{一句话：极限是“最终趋向哪里”，导数是“此刻有多陡”。把这俩搞懂，四章就通了。}

\section{第 1 章 集合与映射：先把话说明白}

\subsection{集合：一袋东西}

\textbf{集合}就是“一袋确定的东西”。比如：
\begin{itemize}[leftmargin=1.4em]
\item $A=\{1,2,3\}$：袋子里装着 1、2、3。
\item “太阳系的行星”是一个集合（装着 8 个行星）。
\item “1 到 10 之间的偶数”= $\{2,4,6,8,10\}$。
\end{itemize}

符号 $\in$ 读作“属于”：$2\in\{1,2,3\}$ 意思是“2 在这个袋子里”。反过来 $4\notin\{1,2,3\}$。

\textbf{子集}：如果甲袋子的每样东西都在乙袋子里，就说甲是乙的子集，写 $A\subseteq B$。比喻：小篮子的东西全装在大篮子里，小篮子就是大篮子的子集。

\textbf{一个有意思的问题：}一个装着 $n$ 样东西的袋子，能分出多少种“小袋子”（子集）？
答案是 $2^n$ 个。为什么？因为袋子里每样东西你都有两种选择：“要”或者“不要”。$n$ 样东西，每样 2 种选择，乘起来就是 $\underbrace{2\times2\times\cdots\times2}_{n\text{ 个}}=2^n$。
（比如 $\{1,2\}$：子集有 $\varnothing,\{1\},\{2\},\{1,2\}$，共 $2^2=4$ 个。$\varnothing$ 是空袋子。）

\idea{集合 = 一袋东西；$n$ 样东西的子集有 $2^n$ 个（每样东西“要/不要”各两种选法）。}

\subsection{映射与函数：给每样东西指定一个“对象”}

\textbf{映射}就是“规定”：左边的每一样东西，都被指定到右边一个东西。

比喻：班里每个学生（左边），都被指定坐到某个座位（右边）。这就是一个“映射”。

\textbf{函数}是映射里最常见的一种——左边是数字，右边也是数字。比如 $y=x^2$：你给它一个 $x$，它吐出 $x^2$。

\textbf{三个必懂的小词：}
\begin{itemize}[leftmargin=1.4em]
\item \textbf{定义域}：允许输入的那些 $x$。比如 $y=\dfrac1x$，$x$ 不能是 0（除以 0 没意义），所以定义域是“除 0 以外的所有数”。
\item \textbf{值域}：函数能吐出的那些 $y$。比如 $y=x^2$ 只能吐出“非负数”（平方不会是负的），值域是 $[0,+\infty)$。
\item \textbf{对应法则}：就是“怎么算”，$y=x^2$ 这个 $x^2$ 就是法则。
\end{itemize}

\textbf{怎么求定义域？}记住三条“禁区”：
\begin{enumerate}[leftmargin=1.4em]
\item \textbf{分母不能为 0}。$\dfrac1{x-2}$，$x\ne2$。
\item \textbf{偶次根号里要 $\ge0$}。$\sqrt{x-3}$，$x\ge3$。
\item \textbf{对数真数要 $>0$}。$\ln(x+1)$，$x>-1$。
\end{enumerate}

\warn{两个函数算“同一个”，必须定义域一样、算法也一样。$y=x$ 和 $y=\dfrac{x^2}{x}$ 看着像，但后者 $x=0$ 没定义，所以不是同一个函数。}

\idea{映射=给左边每个东西指定右边一个；定义域=能输入啥，值域=能吐出啥。求定义域盯住：分母 $\ne0$、根号内 $\ge0$、对数内>0。}

\subsection{复合函数：套娃}

\textbf{复合}就是“把一个函数塞进另一个函数”。比如先算 $u=x^2-3$，再把 $u$ 塞进 $\ln$：得到 $y=\ln(x^2-3)$。

比喻：两条流水线接起来，第一条的产出当第二条的原料。

\textbf{做代换题（很重要，考试常考）：}已知 $f(x+3)=2x+1$，求 $f(x)$。
想法：题目里 $f$ 后面跟着 $x+3$，我想让它跟着“一个字母”。那就\textbf{令 $t=x+3$}，则 $x=t-3$，把 $x$ 全换成 $t-3$：
$$f(t)=2(t-3)+1=2t-5.$$
再把字母 $t$ 改成 $x$：$\boxed{f(x)=2x-5}$。就这么简单——\textbf{令新字母 = 括号里的东西，反解出 $x$，代回去}。

\section{第 2 章 数列极限：一队数最后跑到哪}

\subsection{数列：排成一队的数}

\textbf{数列}就是一串按顺序排的数：$x_1,x_2,x_3,\dots$。比如：
$$1,\ \tfrac12,\ \tfrac13,\ \tfrac14,\ \dots\quad(\text{即 }x_n=\tfrac1n)$$
这一队数越往后越靠近 0。我们就说它的\textbf{极限是 0}，写 $\lim\limits_{n\to\infty}x_n=0$。

\idea{极限 = 这队数“最终无限靠近”的那个数。是“靠近”，不一定是“到达”。}

\subsection{$\eps$–$N$ 语言：把“无限靠近”说成人话}

“无限靠近 0”太含糊，数学家要精确。他们这样想：

你说它“最终跟 0 很近”。那我给你一个“容许误差”$\eps$（念“艾普西龙”），比如 $\eps=0.01$。你要能回答：\textbf{从第几项开始，后面所有项都落进 0 的 $\eps$ 范围里？}

对 $x_n=\frac1n$：要 $\frac1n<0.01$，就是 $n>100$。所以从第 101 项起，后面全都在 $0.01$ 以内。$\eps=0.001$？那就从第 1001 项起。\textbf{无论 $\eps$ 多小，我总能找到“从那以后”的起点。}这就是极限的精确定义：

$$\lim_{n\to\infty}x_n=a\iff\forall\eps>0,\ \exists N,\ \forall n>N:\ |x_n-a|<\eps.$$

\textbf{逐字翻译成人话：}
\begin{itemize}[leftmargin=1.4em]
\item $\forall\eps>0$：随便你给多小的容许误差。
\item $\exists N$：我都能找到一个“起跑线”$N$。
\item $\forall n>N$：起跑线之后的每一项。
\item $|x_n-a|<\eps$：都跟目标 $a$ 差不到 $\eps$。
\end{itemize}

\idea{极限定义就是一场“抬杠游戏”：你出误差 $\eps$，我总能报出一个起点 $N$，保证之后所有项都在误差内。我永远能赢，就是极限存在。}

\warn{顺序绝不能乱！$\forall\eps\,\exists N$ 是“先给你误差，我再找起点”；如果反过来就变成“一个固定的起点对所有误差都管用”，那是错的。$N$ 可以依赖 $\eps$，但不能依赖 $n$（不能说“第 $n$ 项用第 $n$ 项的标准”）。}

\subsection{夹逼定理：两边一夹就定了}

想象三明治：上层面包和下层面包都往 1 靠，中间的火腿还能跑哪去？只能也往 1 靠。

\textbf{夹逼定理：}如果 $a_n\le c_n\le b_n$，而且上下两队都趋于 $L$，那中间 $c_n$ 也一定趋于 $L$。

例子：求 $\lim\left(\frac1{n^2+1}+\frac2{n^2+1}+\cdots+\frac{n}{n^2+1}\right)$。
这个和很难直接算，但可以夹：每一项都 $\ge$ 最小项、$\le$ 最大项，所以
$$\underbrace{\frac{1+2+\cdots+n}{n^2+1}}_{\text{全取最小？不，估算}} \le \text{原式}\le \frac{1+2+\cdots+n}{n^2+1}.$$
（这里分子都是 $1+2+\cdots+n=\frac{n(n+1)}2$，分母都是 $n^2+1$，上下几乎相等。）上下都趋近 $\frac12$，所以原极限 $=\frac12$。

\idea{夹逼 = 三明治。上下界都奔同一个数，中间被逼着也奔那儿。}

\subsection{单调有界：只涨且涨不过头，必然停下}

\textbf{单调}：一队数只会“一直涨”（或一直跌），不会忽上忽下。
\textbf{有界}：一直涨也涨不过某个“天花板”。

\textbf{单调有界定理：}如果一队数一直涨、又涨不过天花板，那它一定会“稳定下来”，趋向某个极限（就是最小的那个天花板，叫上确界）。

比喻：你往杯子里倒水，水位只会涨（单调），但杯子有容量上限（有界），最后水面停在某个高度——不会无限涨下去。

例子：$x_1=\sqrt2$，之后每一项 $x_{n+1}=\sqrt{2+x_n}$。
可以证明它一直涨、又始终小于 2。所以有极限。设极限为 $L$（稳定后 $x_{n+1}$ 和 $x_n$ 都等于 $L$），代入：
$$L=\sqrt{2+L}\Rightarrow L^2=2+L\Rightarrow L^2-L-2=0\Rightarrow L=2\ (\text{或}-1,\text{舍去负数}).$$
所以极限是 2。\textbf{套路：先证有界，再证单调，最后“设极限代入方程解出来”。}

\idea{递推数列求极限三步走：(1) 证明有界 (2) 证明单调 (3) 设 $L$ 代入递推方程，解出 $L$（舍去不合理根）。}

\subsection{确界：天花板的天花板}

\textbf{上确界}（sup）：一堆数里“最小的那个上界”，也就是“最低的天花板”。
比如 $(0,1)$ 里的数，天花板可以是 1、也可以是 2；最低的天花板是 1，所以上确界是 1。

\textbf{关键事实（实数完备性）：}只要一堆数不全无限往上跑（有上界），就\textbf{一定}有一个最低天花板。这叫“确界原理”，是实数比有理数“更完整”的地方——有理数会漏，实数不漏。

\warn{有理数是不够的：$\sqrt2$ 不是有理数，你在一堆“平方小于2的有理数”里找最低天花板，结果那个位置是个“洞”（$\sqrt2$），有理数填不上。实数把洞补上了。}

\section{第 3 章 函数极限与连续：曲线跑到哪}

\subsection{函数极限：把数列的想法搬到曲线上}

数列是“一个一个跳”（$n=1,2,3,\dots$）；函数是“连续滑”（$x$ 可以是 $1.5,1.9,1.999\dots$）。但思想一模一样：

$$\lim_{x\to x_0}f(x)=A$$
意思是：当 $x$ \textbf{无限靠近} $x_0$ 时，$f(x)$ 无限靠近 $A$。

\textbf{$\eps$--$\delta$ 语言}（把“靠近”说精确）：你给误差 $\eps$，我找一个距离 $\delta$（念“德尔塔”），保证“只要 $x$ 离 $x_0$ 不到 $\delta$，$f(x)$ 就离 $A$ 不到 $\eps$”：
$$\forall\eps>0,\ \exists\delta>0:\ 0<|x-x_0|<\delta\Rightarrow|f(x)-A|<\eps.$$

\idea{数列是“你给 $\eps$，我给第 $N$ 项”；函数是“你给 $\eps$，我给一个距离 $\delta$”。换个马甲，还是那场游戏。}

\warn{注意那个 $0<|x-x_0|$：意思是“$x$ 靠近 $x_0$，但不等于 $x_0$”。为什么？“极限”关心的是“靠近时”的趋势，跟 $x_0$ 这一点本身（哪怕函数在那里没定义）没关系。}

\subsection{连续：曲线不断带}

\textbf{连续}就是“曲线画的时候不用抬笔，没有断口”。精确定义：$x$ 靠近 $x_0$ 时，$f(x)$ 也靠近 $f(x_0)$——也就是“极限恰好等于这一点的函数值”。

$$\lim_{x\to x_0}f(x)=f(x_0).$$

对比一下：连续和极限的区别就一句话——\textbf{极限不要求等于 $f(x_0)$，连续要求等于}。

\textbf{间断点}就是“断口”，分三类：
\begin{itemize}[leftmargin=1.4em]
\item \textbf{可去间断}：左右都靠到同一个值，只是那一点“戳了个洞”（函数值和别人不一样，或干脆没定义）。
\item \textbf{跳跃间断}：从左边和从右边靠，落到\textbf{不同}高度，像台阶（跳了一下）。
\item \textbf{无穷间断}：靠近时函数值冲到 $\pm\infty$（像 $\frac1x$ 在 0 附近，冲到无穷）。
\end{itemize}

比喻：画画——可去是“擦了个小点”，跳跃是“画了两段没接上”，无穷是“笔尖冲出了纸”。

\subsection{闭区间上的连续函数：三条金规律}

“闭区间 $[a,b]$”就是“从 $a$ 到 $b$ 一段，两头都算上”。如果函数在这段上连续，有三条好用的结论：

\begin{enumerate}[leftmargin=1.4em]
\item \textbf{有最大值最小值}：这段上一定有个最高的点、最低的点。（画得连续的曲线必有最高峰和最低谷。）
\item \textbf{介值定理}：如果你从高度 1 走到高度 5，中间你一定\textbf{经过}了高度 3。（连续的东西不能“隔空跳”。）
\item \textbf{一致连续（Heine 定理）}：在一段有限的闭区间上，连续的函数“变化速度有统一的控制”——不会出现“某个地方突然陡到没法管”的情况。
\end{enumerate}

\textbf{介值定理的用法（零点定理）：}如果 $f(a)<0$、$f(b)>0$（一负一正），那中间一定有个点让 $f=0$（穿过水平轴）。这个用来证明“方程有解”。

例子：证明 $x=\cos x$ 有解。令 $f(x)=x-\cos x$。$f(0)=0-1=-1<0$，$f(1)=1-\cos1>0$。一负一正，连续，所以中间有个点让 $f=0$，也就是 $x=\cos x$。

\idea{闭区间连续三大宝：(1) 必有最大最小 (2) 中间值都能取到 (3) 一致连续。证方程有根就用“一正一负→零点”。}

\section{第 4 章 微分：这一刻有多陡}

\subsection{导数：瞬时速度}

开车，你从 $t_1$ 到 $t_2$ 走了距离 $\Delta s$，那\textbf{平均速度} = $\frac{\Delta s}{\Delta t}$。但这是“一段路”的平均。如果你想知道“在某一瞬间”的速度呢？让 $\Delta t$ 越来越小，越来越小……

\textbf{导数就是这个“越来越小时”的极限。}写成：
$$f'(x_0)=\lim_{\Delta x\to0}\frac{f(x_0+\Delta x)-f(x_0)}{\Delta x}.$$

\textbf{人话：}在 $x_0$ 这一点，曲线有多陡。几何上它是\textbf{切线斜率}；物理上它是\textbf{瞬时速度}。

\idea{导数 = 曲线在这一点的斜率 = 瞬时速度 = 变化率。就是“$\Delta x$ 缩到 0 时的那个比值的极限”。}

\textbf{微分：}导数乘上小变化量，就是“变化量的近似”。$dy=f'(x)\,dx$。它可以用来做近似估算，比如求“球镀一层铜要多少铜”这种。

\subsection{求导就是背公式 + 用三套工具}

求导大部分是“背公式 + 套规则”。先背几个最基本的：

\begin{center}
\begin{tabular}{c|c}
函数 & 导数\\\hline
$x^n$ & $n x^{n-1}$\\
$\sin x$ & $\cos x$\\
$\cos x$ & $-\sin x$\\
$e^x$ & $e^x$\\
$\ln x$ & $\dfrac1x$\\
常数 $C$ & $0$
\end{tabular}
\end{center}

\textbf{三条规则（求复杂函数靠它们）：}
\begin{enumerate}[leftmargin=1.4em]
\item \textbf{加减}：$(f\pm g)'=f'\pm g'$。（分开求）
\item \textbf{乘积}：$(fg)'=f'g+fg'$。（“前导后不导 + 前不导后导”）
\item \textbf{复合（链式）}：外层导数 $\times$ 内层导数。这是最重要的！$(\sin(x^2))'=\cos(x^2)\cdot 2x$（外层 $\sin$ 求导得 $\cos$，内层 $x^2$ 求导得 $2x$，乘起来）。
\end{enumerate}

\warn{链式法则最容易忘内层！求 $\sin(x^2)$ 的导数时，千万别只写 $\cos(x^2)$，还要乘内层 $x^2$ 的导数 $2x$。想象剥洋葱：一层一层求导，每剥一层就把那层的导数留下。}

\textbf{求隐函数的导数（$y$ 藏在方程里）：}比如 $x^2+y^2=1$，$y$ 不是显式写出来的。那就\textbf{两边同时求导}，遇到 $y$ 记得乘 $y'$：$2x+2yy'=0\Rightarrow y'=-\frac{x}{y}$。

\textbf{参数方程求导：}如果 $x$ 和 $y$ 都由第三个变量 $t$ 表示，那么 $\dfrac{dy}{dx}=\dfrac{dy/dt}{dx/dt}$（分别对 $t$ 求导，再相除）。

\subsection{导数的用处：找最高最低、判增减}

导数正负告诉我们曲线在“涨”还是在“跌”：
\begin{itemize}[leftmargin=1.4em]
\item $f'>0$：曲线在涨（上升）。
\item $f'<0$：曲线在跌（下降）。
\item $f'=0$：曲线在这一点“水平”，可能是最高点、最低点（也可能不是，叫“驻点”）。
\end{itemize}

所以找最值：先找 $f'=0$ 的点，再比较它们和两端点的函数值，最大的就是最大值。

\idea{求最值：(1) 求导 (2) 令导数为 0 找到“可疑点” (3) 把可疑点和端点都代进去比大小。}

\section{考前最后嘱咐}

\begin{enumerate}[leftmargin=1.4em]
\item \textbf{每章一句话：}
  \begin{itemize}
  \item 第1章：集合=一袋东西，映射=指定对应，定义域盯三个禁区。
  \item 第2章：极限=最终靠近哪；$\eps$–$N$=“你给误差我找起点”的游戏；夹逼=三明治；单调有界=只涨不过头必停。
  \item 第3章：连续=不断带；闭区间三宝：最值、介值、一致连续；证方程有根=一正一负。
  \item 第4章：导数=斜率/速度；求导=背公式+链式；最值=导数为0处比较。
  \end{itemize}
\item \textbf{动手比看更重要：}光看不行，把你做错的题抄下来，隔天重做一遍。
\item \textbf{看不懂公式就先想比喻：}“靠近”“三明治”“只涨不过头”“斜率”这几个画面抓住，公式只是把它们写严谨。
\end{enumerate}

\end{document}
